% Use IEEE Transaction format (2-cols, standard for AI/CS papers)
\documentclass[10pt, journal, compsoc]{IEEEtran}

% Import our customized styling module
\usepackage{setting/khang_paper}

\begin{document}

\title{Gaussian Mixture Model Clustering with Incomplete Data:\\
    From Paper Reproduction to A Research-Oriented Technical Report}

\author{
    H.Khang Nguyen , T.Nhat Nguyen and L.Phuc Nguyen
    \IEEEcompsocitemizethanks{\IEEEcompsocthanksitem All authors are with the Master of Data Science, University of Science, VNU-HCM, Ho Chi Minh City, Vietnam.
        \protect\\
        E-mail: \{25C0103580, \}@student.hcmus.edu.vn}
}

\IEEEtitleabstractindextext{%
    \begin{abstract}
        This report revisits the paper ``Gaussian Mixture Model Clustering with Incomplete Data'' from a research-oriented perspective rather than a plain recap. We reconstruct the underlying probabilistic model with a consistent notation system, reformulate the method as a constrained maximum-likelihood problem over both Gaussian mixture parameters and latent missing entries, and clarify the role of the alternating expectation-maximization updates. Beyond restating the original method, we analyze why dynamic imputation can outperform two-stage pipelines, explain the closed-form missing-data update through precision-matrix block structure, discuss convergence guarantees and their limits, and identify practical weaknesses related to initialization, covariance regularization, missingness assumptions, and model selection. The final outcome is a modular \LaTeX{} report intended as a stronger technical foundation for implementation, extension, and experimental redesign.
    \end{abstract}

    \begin{IEEEkeywords}
        Gaussian Mixture Model, Incomplete Data, Missing Data Imputation, EM Algorithm, Model-Based Clustering.
    \end{IEEEkeywords}}

\maketitle
\IEEEdisplaynontitleabstractindextext
\IEEEpeerreviewmaketitle

\section{Introduction}
Clustering under incomplete observations is a more realistic setting than the fully observed scenario assumed by most classical mixture models. In many applied pipelines, missing entries appear because of sensor failure, costly measurements, manual annotation gaps, or heterogeneous acquisition protocols. A standard workaround is to first impute the missing values and then apply a clustering algorithm to the completed matrix. While operationally simple, that strategy disconnects the imputation objective from the clustering objective. The filled values may look statistically plausible at the feature level, yet still be suboptimal for separating latent groups.

The paper by Zhang \textit{et al.}~\cite{zhang2021gmm} addresses this mismatch by coupling imputation and clustering in a single iterative procedure. Their main idea is concise: the current Gaussian mixture model (GMM) parameters induce a refined estimate of the missing entries, and the newly completed data matrix is then used to update the GMM again. Instead of treating imputation as a preprocessing routine, the method promotes it to an optimization variable.

This report keeps the core algorithmic content of the original paper, but its goal is different. We do not merely summarize the article. We rebuild the method with cleaner notation, expose the statistical meaning of each update, and emphasize what the paper leaves implicit:
\begin{itemize}
    \item the optimization is a constrained maximum-likelihood problem over both model parameters and missing values;
    \item the data-completion step is a responsibility-weighted conditional reconstruction in the precision domain;
    \item the convergence argument is local and monotonic, not a guarantee of global optimality;
    \item the empirical gains are strongest when classical two-stage imputers fail to preserve cluster structure.
\end{itemize}

Our perspective is that the method is best understood as a \highlight{one-stage model-based clustering framework for incomplete data}, with direct connections to alternating optimization, latent-variable inference, and conditional Gaussian estimation. This viewpoint also makes it easier to see where the method can be improved, for example by using structured covariance models, soft missingness mechanisms, or automatic model selection for the number of clusters.


\section{Notation and Problem Setup}
We adopt a notation system that strictly distinguishes scalars, vectors, and matrices:
\begin{itemize}
    \item scalars use plain math symbols, for example $\scal{n}$, $\scal{d}$, $\scal{K}$, and $\scal{\alpha}_k$;
    \item vectors use $\vect{\cdot}$, for example $\vect{x}_j \in \R^{\scal{d}}$ and $\vect{\mu}_k \in \R^{\scal{d}}$;
    \item matrices use $\mat{\cdot}$, for example $\mat{X} \in \R^{\scal{n}\times \scal{d}}$ and $\mat{\Sigma}_k \in \R^{\scal{d}\times \scal{d}}$.
\end{itemize}

\begin{table}[t]
    \caption{Core notation used throughout the report}
    \label{tab:notation}
    \centering
    \begin{tabularx}{\linewidth}{@{}lX@{}}
        \toprule
        Symbol & Meaning \\
        \midrule
        $\mat{X} \in \R^{\scal{n}\times \scal{d}}$ & Completed data matrix whose rows are samples \\
        $\vect{x}_j$ & The $\scal{j}$th sample vector \\
        $\vect{x}_{j,o}, \vect{x}_{j,m}$ & Observed and missing blocks of sample $\vect{x}_j$ \\
        $\scal{K}$ & Number of Gaussian mixture components \\
        $\scal{\alpha}_k$ & Mixture weight of component $\scal{k}$ \\
        $\vect{\mu}_k$ & Mean vector of component $\scal{k}$ \\
        $\mat{\Sigma}_k$ & Covariance matrix of component $\scal{k}$ \\
        $\scal{\gamma}_{jk}$ & Responsibility of component $\scal{k}$ for sample $\scal{j}$ \\
        $\mathcal{O}_j, \mathcal{M}_j$ & Index sets of observed and missing coordinates \\
        \bottomrule
    \end{tabularx}
\end{table}

Let the dataset be
\[
    \mat{X} =
    \begin{bmatrix}
        \vect{x}_1\trans\\
        \vect{x}_2\trans\\
        \vdots\\
        \vect{x}_{\scal{n}}\trans
    \end{bmatrix},
    \qquad
    \vect{x}_j \in \R^{\scal{d}}.
\]
For each sample $\vect{x}_j$, let $\mathcal{O}_j$ be the set of observed feature indices and $\mathcal{M}_j$ be the set of missing feature indices, with $\mathcal{O}_j \cup \mathcal{M}_j = \set{1,\dots,\scal{d}}$ and $\mathcal{O}_j \cap \mathcal{M}_j = \varnothing$. We partition the sample into
\[
    \vect{x}_j =
    \begin{bmatrix}
        \vect{x}_{j,o}\\
        \vect{x}_{j,m}
    \end{bmatrix},
\]
where $\vect{x}_{j,o}$ is fixed by observation and $\vect{x}_{j,m}$ is unknown.

For a classical GMM with $\scal{K}$ components, the density of a fully observed sample is
\[
    p\paren*{\vect{x}_j}
    =
    \sum_{k=1}^{\scal{K}}
    \scal{\alpha}_k
    \N\paren*{\vect{x}_j \given \vect{\mu}_k, \mat{\Sigma}_k},
\]
where $\scal{\alpha}_k \ge 0$ and $\sum_{k=1}^{\scal{K}} \scal{\alpha}_k = 1$.

The original paper reformulates incomplete-data clustering as the following constrained problem:
\begin{equation}
\begin{aligned}
    &\argmax_{\mat{X},\, \set*{\scal{\alpha}_k,\vect{\mu}_k,\mat{\Sigma}_k}_{k=1}^{\scal{K}}}
    \sum_{j=1}^{\scal{n}}
    \log\paren*{
        \sum_{k=1}^{\scal{K}}
        \scal{\alpha}_k
        \N\paren*{\vect{x}_j \given \vect{\mu}_k, \mat{\Sigma}_k}
    } \\
    &\text{s.t. }
    \vect{x}_{j,o} = \vect{x}^{\,\text{obs}}_{j,o},\ \forall j.
\end{aligned}
\end{equation}

This formulation matters conceptually. It says that the completed data matrix is not an external preprocessing artifact; it is part of the optimization landscape. That single modeling choice is what separates this method from mean imputation, zero filling, or even EM-based imputers used as standalone preprocessing~\cite{ghahramani1994em,aste2015incomplete}.

The latent assignment variable $\scal{z}_j \in \set{1,\dots,\scal{K}}$ is still present as in standard GMM inference. Therefore, the problem contains two layers of hidden structure:
\begin{itemize}
    \item \highlight{cluster uncertainty}, represented by posterior responsibilities;
    \item \highlight{feature uncertainty}, represented by the missing subvectors $\vect{x}_{j,m}$.
\end{itemize}

This dual latent structure is exactly why a one-shot estimator is hard to derive directly and why the paper moves to an alternating optimization scheme.

\begin{remark}
The optimization variable is the \emph{completed} matrix $\mat{X}$, not the raw incomplete matrix. This subtle distinction is what turns the method from a preprocessing pipeline into a joint estimation framework.
\end{remark}


\section{Methodology}
The proposed algorithm alternates between estimating cluster memberships, updating GMM parameters, and refining the missing entries. It can be interpreted as an EM-like procedure nested inside a block-coordinate ascent routine.

\subsection{Responsibility Update}
Given the current completed matrix $\mat{X}$ and current parameters $\set*{\scal{\alpha}_k,\vect{\mu}_k,\mat{\Sigma}_k}_{k=1}^{\scal{K}}$, define the posterior responsibility of component $\scal{k}$ for sample $\vect{x}_j$ by
\[
    \scal{\gamma}_{jk}
    =
    p\paren*{\scal{z}_j = \scal{k} \given \vect{x}_j}
    =
    \frac{
        \scal{\alpha}_k \N\paren*{\vect{x}_j \given \vect{\mu}_k, \mat{\Sigma}_k}
    }{
        \sum_{\ell=1}^{\scal{K}}
        \scal{\alpha}_{\ell} \N\paren*{\vect{x}_j \given \vect{\mu}_{\ell}, \mat{\Sigma}_{\ell}}
    }.
\]
This is the standard E-step of model-based clustering~\cite{bouman1997cluster,melnykov2010finite}. The difference from the classical case is that $\vect{x}_j$ itself contains values generated by the previous completion step.

\subsection{Parameter Update with Fixed Completed Data}
Once $\scal{\gamma}_{jk}$ is available, the mixture coefficients, means, and covariance matrices are updated exactly as in a conventional GMM fitted on the current completed matrix:
\[
    \scal{N}_k = \sum_{j=1}^{\scal{n}} \scal{\gamma}_{jk},
    \qquad
    \scal{\alpha}_k = \frac{\scal{N}_k}{\scal{n}},
\]
\[
    \vect{\mu}_k
    =
    \frac{1}{\scal{N}_k}
    \sum_{j=1}^{\scal{n}}
    \scal{\gamma}_{jk}\vect{x}_j,
\]
\[
    \mat{\Sigma}_k
    =
    \frac{1}{\scal{N}_k}
    \sum_{j=1}^{\scal{n}}
    \scal{\gamma}_{jk}
    \paren*{\vect{x}_j - \vect{\mu}_k}
    \paren*{\vect{x}_j - \vect{\mu}_k}\trans.
\]

These equations reveal an important point often glossed over in short summaries: once the missing entries are temporarily filled, the model does not use a special incomplete-data likelihood at this stage. It uses the ordinary complete-data GMM update on the current surrogate matrix. The novelty lies in how that surrogate matrix is refreshed.

\subsection{Missing-Data Update with Fixed GMM Parameters}
The most distinctive step of the method is the update of $\vect{x}_{j,m}$. With all GMM parameters fixed, each sample solves
\[
    \argmax_{\vect{x}_{j,m}}
    \sum_{k=1}^{\scal{K}}
    \scal{\alpha}_k
    \N\paren*{\vect{x}_j \given \vect{\mu}_k, \mat{\Sigma}_k}
    \quad
    \text{s.t. }
    \vect{x}_{j,o} = \vect{x}^{\,\text{obs}}_{j,o}.
\]

To make the derivation explicit, partition each mean vector and precision matrix according to observed and missing coordinates:
\[
    \vect{\mu}_k =
    \begin{bmatrix}
        \vect{\mu}_{k,o}\\
        \vect{\mu}_{k,m}
    \end{bmatrix},
    \qquad
    \mat{\Sigma}_k\inv =
    \begin{bmatrix}
        \mat{\Lambda}_{k,oo} & \mat{\Lambda}_{k,om}\\
        \mat{\Lambda}_{k,mo} & \mat{\Lambda}_{k,mm}
    \end{bmatrix}.
\]
The update derived in the paper is
\begin{equation}
\begin{aligned}
    \vect{x}_{j,m}
    &=
    \paren*{
        \sum_{k=1}^{\scal{K}}
        \scal{P}_{jk}\mat{\Lambda}_{k,mm}
    }\inv \\
    &\quad \times
    \sum_{k=1}^{\scal{K}}
    \scal{P}_{jk}
    \paren*{
        \mat{\Lambda}_{k,mo}\vect{\mu}_{k,o}
        +
        \mat{\Lambda}_{k,mm}\vect{\mu}_{k,m}
        -
        \mat{\Lambda}_{k,mo}\vect{x}_{j,o}
    },
\end{aligned}
\end{equation}
where
\[
    \scal{P}_{jk} = \scal{\alpha}_k \N\paren*{\vect{x}_j \given \vect{\mu}_k, \mat{\Sigma}_k}.
\]

This equation is the mathematical core of the method. It says that the missing block is reconstructed from a weighted combination of component-wise conditional geometry. The weights depend on current mixture plausibility, and the linear system is expressed in the \highlight{precision domain} rather than directly in covariance form. That is why the update resembles a compromise across local quadratic energies induced by each Gaussian component.

\subsection{Interpretation as Conditional Gaussian Reconstruction}
For a single Gaussian component, the optimal missing-value estimate under a quadratic log-density would reduce to a conditional mean expression driven by the relation between observed and missing coordinates. The proposed method generalizes that logic to a mixture model. Instead of selecting one component, it interpolates across components through the weights $\scal{P}_{jk}$. This yields a softer and more adaptive estimator than hard assignment to a single cluster center.

\begin{table}[t]
    \caption{Two-stage imputation versus the one-stage GMM framework}
    \label{tab:one-stage-vs-two-stage}
    \centering
    \begin{tabularx}{\linewidth}{@{}lX X@{}}
        \toprule
        Aspect & Two-stage pipeline & Proposed one-stage method \\
        \midrule
        Imputation target & Generic feature reconstruction & Cluster-aware likelihood improvement \\
        Feedback from clustering & None & Directly affects missing-value updates \\
        Use of covariance geometry & Often weak or absent & Explicit through $\mat{\Sigma}_k$ and $\mat{\Sigma}_k\inv$ \\
        Robustness at high missingness & Tends to deteriorate quickly & Usually more stable if initialization is adequate \\
        Optimization view & Sequential heuristics & Alternating ascent on a joint objective \\
        \bottomrule
    \end{tabularx}
\end{table}

\subsection{Algorithmic Interpretation}
The whole procedure can be summarized as follows:
\begin{enumerate}
    \item initialize missing entries with a simple imputer such as column means;
    \item initialize GMM parameters on the completed matrix;
    \item compute posterior responsibilities $\scal{\gamma}_{jk}$;
    \item update $\scal{\alpha}_k$, $\vect{\mu}_k$, and $\mat{\Sigma}_k$;
    \item update each $\vect{x}_{j,m}$ while keeping $\vect{x}_{j,o}$ unchanged;
    \item repeat until the log-likelihood improvement falls below a tolerance.
\end{enumerate}

This is best viewed as block-coordinate ascent on a non-convex objective. Each block update is designed to avoid decreasing the current objective, so the procedure inherits a monotonic improvement property. However, because the optimization is non-convex, the destination is a local optimum rather than a globally optimal solution.

\subsection{Complexity and Numerical Stability}
The paper reports a complexity of $\mathcal{O}\paren*{\scal{t}\scal{K}\scal{n}\scal{d}^3}$, with $\scal{t}$ iterations, $\scal{K}$ clusters, $\scal{n}$ samples, and $\scal{d}$ dimensions. The cubic term comes from matrix inversion in the data-completion step. In practice, this immediately implies three engineering concerns:
\begin{itemize}
    \item covariance regularization is mandatory when clusters are small or features are highly correlated;
    \item poor initialization can trap the method in unstable local optima;
    \item the method becomes expensive in high dimensions unless one constrains $\mat{\Sigma}_k$, for example to diagonal, tied, or low-rank structures.
\end{itemize}

These issues are not implementation details; they shape the actual usability of the method.

\subsection{Theoretical Boundary of the Current Formulation}
The method remains elegant, but it is not equivalent to exact maximum likelihood under incomplete observations in the classical missing-data EM sense. A textbook incomplete-data EM would integrate over missing entries in the E-step, whereas this method updates them by deterministic optimization. That difference is not a flaw by itself, but it must be stated clearly because it affects how we interpret the algorithm's statistical guarantees.


\section{Implementation and Reproducibility}
This project already contains a working Python implementation in the file \texttt{src/gmm\_missing.py}, which makes it possible to connect the mathematical derivation to concrete engineering choices. That link is valuable because several numerically important details do not appear explicitly in the original paper.

\subsection{Algorithm Box}
\begin{figure}[t]
\hrule
\vspace{0.5em}
\begin{algorithmic}[1]
\Require incomplete matrix $\mat{X}^{\text{inc}}$, number of clusters $\scal{K}$, tolerance $\scal{\varepsilon}$, maximum iterations $\scal{T}_{\max}$
\Ensure completed matrix $\mat{X}$, responsibilities $\mat{\Gamma}$, labels $\vect{\lambda}$
\State initialize $\mat{X}$ by column-mean imputation
\State initialize $\set*{\scal{\alpha}_k,\vect{\mu}_k,\mat{\Sigma}_k}_{k=1}^{\scal{K}}$
\For{$\scal{t}=1,\dots,\scal{T}_{\max}$}
    \State compute responsibilities $\scal{\gamma}_{jk}$ for all $(j,k)$
    \State update $\scal{\alpha}_k$, $\vect{\mu}_k$, and $\mat{\Sigma}_k$
    \For{each sample $\vect{x}_j$ with missing block $\vect{x}_{j,m}$}
        \State update $\vect{x}_{j,m}$ by the precision-weighted closed form
        \State keep $\vect{x}_{j,o}$ unchanged
    \EndFor
    \If{relative log-likelihood improvement $\le \scal{\varepsilon}$}
    \State \textbf{break}
    \EndIf
\EndFor
\State $\lambda_j \gets \argmax_k\ \scal{\gamma}_{jk}$ for each sample
\end{algorithmic}
\vspace{0.5em}
\hrule
\caption{One-stage GMM clustering with incomplete data.}
\label{alg:gmm-missing}
\end{figure}

\subsection{Numerical Safeguards}
The code-level implementation introduces several safeguards that a professional report should document explicitly:
\begin{itemize}
    \item \textbf{Mean initialization for missing values.} A simple initial completion is used only to start the alternating routine; it is not the final imputation mechanism.
    \item \textbf{K-means parameter initialization.} Component means are initialized from K-means centers, which generally improves stability over random initialization.
    \item \textbf{Covariance regularization.} A diagonal ridge term is added to each covariance estimate to prevent singularity when some clusters are small or nearly collinear.
    \item \textbf{Log-domain responsibility computation.} Using log-likelihoods and a log-sum-exp normalization is more stable than evaluating Gaussian densities directly.
    \item \textbf{Fallback for singular precision blocks.} If a block inversion fails during the update of $\vect{x}_{j,m}$, the implementation retains the current value rather than crashing.
\end{itemize}

\subsection{Mapping the Equations to the Code}
The implementation mirrors the paper almost line by line. The E-step computes a matrix $\mat{\Gamma} \in \R^{\scal{n}\times \scal{K}}$ whose entries are the responsibilities $\scal{\gamma}_{jk}$. The parameter M-step then forms
\[
    \scal{N}_k = \sum_{j=1}^{\scal{n}} \scal{\gamma}_{jk},
    \qquad
    \vect{\mu}_k = \frac{1}{\scal{N}_k}\sum_{j=1}^{\scal{n}}\scal{\gamma}_{jk}\vect{x}_j,
\]
and
\[
    \mat{\Sigma}_k
    =
    \frac{1}{\scal{N}_k}
    \sum_{j=1}^{\scal{n}}
    \scal{\gamma}_{jk}
    \paren*{\vect{x}_j-\vect{\mu}_k}
    \paren*{\vect{x}_j-\vect{\mu}_k}\trans
    + \scal{\lambda}\mat{I},
\]
where $\scal{\lambda}>0$ is the regularization level used in code.

The data-update step then loops over samples, extracts missing and observed coordinate sets, partitions the component precision matrices accordingly, accumulates the weighted system matrix and right-hand side, and solves for $\vect{x}_{j,m}$. This is exactly the operational form of the paper's Equation (14).

\subsection{Reproducibility Recommendations}
To make the report more reproducible than the source paper, future experiments should log at least the following metadata:
\begin{itemize}
    \item random seed for missing-pattern generation and model initialization;
    \item covariance regularization level and stopping tolerance;
    \item average runtime and iteration count for each dataset and missing ratio;
    \item failure counts due to singular updates or degenerate clusters;
    \item sensitivity curves with respect to $\scal{K}$, initialization, and missingness severity.
\end{itemize}

\begin{remark}
If these metadata are omitted, a method may appear empirically strong while actually depending on a narrow band of stable hyperparameter settings. That risk is especially serious for non-convex mixture models.
\end{remark}


\section{Experimental Re-Reading}
The original paper evaluates the method on eight UCI datasets: Iris, AlcoholQCM, Seeds, Wine, Segment, ElectricalGrid, Avila, and Letter~\cite{zhang2021gmm}. Missing entries are generated synthetically from fully observed data at ratios from $10\%$ to $70\%$. The algorithm is compared against three static imputers, namely mean filling, zero filling, and EM filling, as well as dynamic K-means filling variants from earlier work~\cite{wang2019kmeans}.

\begin{table}[t]
    \caption{A stronger experimental protocol for a modern rewrite}
    \label{tab:protocol}
    \centering
    \begin{tabularx}{\linewidth}{@{}lX@{}}
        \toprule
        Item & Recommendation \\
        \midrule
        Missingness regimes & Evaluate MCAR, MAR, and stress-test MNAR-inspired corruption \\
        Cluster count & Report results with true $\scal{K}$ and model-selected $\scal{K}$ \\
        Metrics & ACC, NMI, F-score, purity, runtime, iterations, failure rate \\
        Stability & Repeat across multiple seeds and missing-pattern draws \\
        Ablations & Initialization, covariance type, regularization, stopping tolerance \\
        Baselines & Include both classical and modern incomplete-data clustering methods \\
        \bottomrule
    \end{tabularx}
\end{table}

\subsection{What the Reported Results Actually Show}
Across ACC, NMI, F-score, and purity, the paper consistently reports that the proposed method outperforms classical two-stage pipelines. The strongest pattern is not simply ``GMM is better than K-means.'' A more precise reading is:
\begin{itemize}
    \item when imputation is decoupled from clustering, the filled matrix can erase discriminative structure;
    \item when imputation is updated in tandem with cluster assignments, the filled values become more cluster-aware;
    \item the advantage grows as the missing ratio becomes moderate or high, where naive imputers collapse faster.
\end{itemize}

This interpretation is more informative than listing performance tables. It connects the empirical results back to the modeling decision introduced in the formulation.

\subsection{Local Reproduction on \texttt{CC GENERAL}}
Our project extends beyond the original paper by reproducing the method on a separate real tabular dataset, \texttt{CC GENERAL}, through the notebook \texttt{GMM\_Nhat.ipynb}. This dataset contains $8950$ credit-card customers and $17$ numeric behavioral features after dropping the customer identifier. In the stored notebook run, we inject approximately $30\%$ missingness, which yields $45{,}789$ missing entries.

The reproduction notebook then walks through the same logic as the paper:
\begin{itemize}
    \item initialize missing values by mean imputation;
    \item standardize the completed matrix;
    \item initialize GMM parameters with K-means;
    \item compute responsibilities and perform one explicit EM-style update;
    \item run the full alternating routine with \texttt{fit\_scale()}.
\end{itemize}

The recorded run converges after $199$ iterations and yields a silhouette score of $0.1489$. We do not claim that this value is intrinsically high; its importance is different. It demonstrates that the model can be trained end-to-end on a moderate-dimensional financial dataset with large injected missingness while still producing a coherent partition structure.

\begin{figure}[t]
    \centering
    \includegraphics[width=\linewidth]{notebook_exports/GMM_Nhat_cell8_out1.png}
    \caption{Reproduction artifact from \texttt{GMM\_Nhat.ipynb}: clustered samples after dynamic imputation, with previously incomplete rows highlighted. This figure is useful because it exposes where imputed samples land in the learned geometry.}
    \label{fig:cc-general-reproduction}
\end{figure}

\subsection{Why Dynamic Imputation Helps}
Suppose two clusters overlap in marginal feature means but differ in multivariate shape. A static imputer based on global averages or independent feature heuristics may pull incomplete samples toward a global center. In contrast, the one-stage GMM method uses posterior responsibilities and covariance structure to complete missing values in a way that remains compatible with latent cluster geometry. In that sense, the algorithm does not merely repair missingness; it \highlight{preserves separability}.

\subsection{Limits of the Experimental Design}
Even though the empirical section is convincing at first glance, several limitations should be acknowledged explicitly in a stronger report:
\begin{itemize}
    \item the missingness is synthetically generated from complete datasets, which primarily matches MCAR-style setups;
    \item the number of clusters is assumed known in advance, removing an important practical difficulty;
    \item there is no ablation on initialization, covariance regularization, or sensitivity to nearly singular clusters;
    \item the baselines do not include more recent incomplete-data clustering methods or modern probabilistic deep approaches;
    \item performance gains are shown, but runtime and memory tradeoffs are not analyzed with the same care.
\end{itemize}

\subsection{What a Better Reproduction Should Add}
A stronger reproduction should go beyond reproducing aggregate scores. It should also reveal \emph{why} the method wins and \emph{when} it fails. In practical terms, the following additions would materially improve the experimental section:
\begin{itemize}
    \item report per-dataset runtime and iteration counts to quantify the price paid for one-stage optimization;
    \item compare full covariance against diagonal and tied covariance variants to isolate the role of geometric flexibility;
    \item measure imputation error together with clustering quality to see whether better clustering always aligns with better reconstruction;
    \item examine sensitivity under class imbalance, since mixture models can behave differently when some components are weakly represented.
\end{itemize}

Our current notebooks already move in that direction because they preserve intermediate outputs, convergence summaries, and visual artifacts. The next step is to convert those artifacts into a systematic experimental log rather than leaving them as notebook-only evidence.

\subsection{A Curious Observation Worth Preserving}
The paper notes that for some datasets, performance with moderate missingness can exceed performance on the complete data. This is counterintuitive, but not impossible. A plausible explanation is that the iterative completion step can smooth noise or attenuate outliers by pulling uncertain coordinates toward cluster-consistent regions. That does not mean missingness is beneficial in general; rather, it suggests the algorithm may implicitly act as a form of regularized denoising. This is one of the most interesting observations in the article and deserves more attention than it receives in the original presentation.


\section{Critical Discussion and Extensions}
The original contribution is solid, but its deeper value appears when we reinterpret it beyond the paper's exact wording.

\subsection{Key Insight 1: The Method Couples Geometry and Completion}
The method is stronger than a pipeline of ``fill then cluster'' because it allows cluster geometry to directly influence the missing entries. This is especially important for multimodal data, where marginal feature averages are poor summaries. The covariance matrices $\mat{\Sigma}_k$ encode local anisotropy, and the posterior weights determine which local geometry matters for each sample. Therefore, the completion step is not generic imputation; it is \highlight{cluster-conditional geometric reconstruction}.

\subsection{Key Insight 2: The Objective Is Only Partially Probabilistic}
Although the algorithm is inspired by likelihood maximization, the completion step optimizes missing coordinates as point estimates rather than integrating them out. In a fully Bayesian treatment, one would marginalize over the missing values or sample from their conditional posterior. The paper instead chooses deterministic optimization. This makes the algorithm simpler and computationally tractable, but it also means uncertainty in the imputed entries is collapsed too early.

\subsection{Key Insight 3: Convergence Is Monotone but Local}
The paper is correct in claiming monotonic improvement under the alternating updates, but that should not be overstated. Because the joint objective is non-convex in $\mat{X}$ and the mixture parameters, the algorithm can converge to poor local maxima, especially under weak initialization or heavily overlapping clusters. A more rigorous phrasing is that the method is guaranteed to produce a non-decreasing objective sequence and converge to a stationary point under standard regularity assumptions.

\subsection{Practical Failure Modes}
For real deployments, the following risks deserve explicit mention:
\begin{itemize}
    \item if too many coordinates are missing in the same region, covariance estimation becomes brittle;
    \item if a component receives very low effective mass $\scal{N}_k$, then $\mat{\Sigma}_k$ may become nearly singular;
    \item if the true missingness mechanism is not MCAR or MAR, the learned clusters may absorb systematic measurement bias;
    \item if $\scal{K}$ is misspecified, the imputation step can become self-reinforcing in the wrong latent partition.
\end{itemize}

\subsection{Extensions That Would Genuinely Improve the Paper}
Several research directions follow naturally from the current framework:
\begin{itemize}
    \item \textbf{Structured covariance models.} Replacing full covariance matrices by diagonal, tied, sparse, or factor-analyzer forms would improve scalability and stability.
    \item \textbf{Automatic model selection.} Integrating BIC, ICL, or cross-validated likelihood would remove the unrealistic assumption that $\scal{K}$ is known.
    \item \textbf{Missingness-aware modeling.} A stronger extension would explicitly model the missingness process, especially when absence is informative rather than random.
    \item \textbf{Uncertainty-preserving completion.} Instead of a single point estimate for $\vect{x}_{j,m}$, one could propagate conditional variance into later updates.
    \item \textbf{Hybrid deep probabilistic models.} The same one-stage idea could be adapted to latent-variable models where a neural encoder parameterizes mixture components in a lower-dimensional space.
\end{itemize}

\subsection{A Useful Conceptual Reframing}
One productive way to reframe the paper is to think of it as solving a \highlight{task-aware imputation} problem. The method does not attempt to estimate missing values in a universally optimal sense. Instead, it estimates them in the way that best supports a downstream clustering objective under a parametric mixture model. This perspective is useful because it generalizes immediately to other unsupervised tasks such as subspace learning, graph clustering, or representation learning with partial observations.

\subsection{How This Report Improves on a Plain Recap}
Compared with a standard paper summary, this report contributes three upgrades. First, it normalizes the mathematical notation so that the role of each object is unambiguous. Second, it explains the closed-form completion step as a weighted precision-based conditional estimator, which is the real intellectual center of the algorithm. Third, it reframes the empirical claims as evidence for a broader principle: \highlight{imputation should be optimized for the downstream unsupervised objective, not for standalone reconstruction quality alone}. That principle is reusable far beyond GMM clustering.


\section{Conclusion}
This report reconstructed the paper ``Gaussian Mixture Model Clustering with Incomplete Data'' as a cleaner and more critical technical document. The main algorithmic idea is to optimize the completed data matrix and the GMM parameters jointly, so that imputation is no longer detached from clustering. Once written in a consistent notation, the method becomes easier to understand: the E-step computes responsibilities, the standard GMM M-step updates mixture parameters, and the distinctive data-completion step solves a precision-weighted conditional reconstruction problem for each incomplete sample.

Beyond reproducing the paper, we emphasized what makes the method scientifically interesting. Its empirical advantage comes from aligning missing-value completion with cluster geometry, not from imputation quality in isolation. At the same time, its guarantees remain local, its computational cost is nontrivial, and its practical reliability depends on initialization, covariance regularization, the missingness mechanism, and the choice of the number of clusters.

As a result, this paper should be read not only as a standalone algorithm, but also as a template for a broader design principle in unsupervised learning with incomplete observations: \highlight{latent structure and data completion should be learned together whenever possible}. That principle opens the door to stronger probabilistic, scalable, and missingness-aware extensions in future work.


\bibliographystyle{IEEEtran}
\bibliography{bib/references}

\end{document}
